% Pursuing Stacks (cotes 134-2 to 134-8): archive cote 134-3, site page 3 (PDF page 4), typescript page 3.
% typescript. Draft compared: ps2.tex:104-159.
% AI-transcribed by Claude Opus 5.5 (claude-opus-5-5, adaptive thinking, effort high), Message Batches API, batch ps1 (msgbatch_01Rv5hCaNTtUt4JHx1yqoHPL), 2026-11-12; not yet checked by eye.
\scanpage{134-3/3}
\margin{3}
\mspage{3}
pro\unsure{c}ess (formally analogyous to the passage from \xout{chain} categories of
chain complexes to the corresponding derived categories), to essentially
``the same'' (up to \unsure{equivalence})\hand{|}``homotopy category''. One way of describing
it is via topological spaces which are not ``too wild'' as objects (the
\unsure{C.W.-spaces}), \deleted{and \xout{maps}}\inserted[typed]{morphisms} being homotopy classes of continuous maps. The
other is via semi-simplicial sets, taking for instance Kan complexes as
objects, and again homotopy classes of ``maps'' as morphisms. The first
description is the one most adapted to direct topological intuition,
as long a\over{s}{t} least as no more adequate notion than ``topological spaces'' is
at hand. The second has the advantage of being a purely algebraic
description, with rather amazing conceptual simplicity moreover. In
terms of the two basic sets of algebraic invariants of a space which hav\unsure{e}
turned up so far, namely cohomology (or homology) on the one hand, and
homotopy gro\hand{u}ps on the other, it can be said that the description via
topological spaces is adquate for direct \deleted{description}\inserted{computation} of neither \inserted{co}homology
nor homotopy groups, whereas the description via semi-simplicial sets i\hand{s}
fairly adequate for \deleted{description}\inserted{computation} of cohomology groups (taking simply the
ab\unsure{e}lianization of the semi-simplicial set, which turns out to be a
chain complex, and taking it's homology and cohomology groups). The
same can be said for the alternative algebraic description of homotopy
types, using cubical complexes instead of semi-simplicial ones, which
were introduced by Serre\hand{,} as they were better suited, it seems,\hand{(2),} for the
study of fibrations and \inserted{of} the homology and cohomology spectral sequences
relative to these. One somewhat surprising common feature of those two
standard algebraic descriptions of homotopy types, is that neither is
any better adapted for a direct desc\hand{r}iption of homotopy groups than the
objects we started with, namely topological spaces. This is all the
more remarkable as it is the homotopy groups really, rather than the
cohomology groups, which are commonly viewed as \uline{the} basic inveriants
in the homotopy point of view, sufficient e.g. for testing whether a
given map \xout{given} is a ``weak equi\uline{valence''} namely gives rise to an isomor-%
phism in the homotopy \uline{category.}\hand{||}\margin{$\lrcorner$||} It is here of course that the point of
view of ``stacks'' (``champs'' in French) of (L) (previously called ``$\infty$-\hand{grou-}
\unsure{poïdes|poïds}'' in the beginnings of the reflections of (L)) sets in. These
presumably give rise to a ``category of models'' and from there, to the
usual homotopy category, in much the same way as topological spaces
or simplicial (or cubical) complexes, thus yielding a third main des-%
cription\inserted{of homotopy types,}, and corresponding wealth of algebraico-geometric intuitions\over{,}{.}
\deleted{of homotopy types,} Moreover, stacks are ideally suited for expressing
the homotopy groups, in an even more direct way than simplicial comple-%
